% 20-09(20쪽) — y=f(x) 의 그래프. x<0 은 기울기 -1 로 원점에 닿고, f(0)=0(찬 점), x>0 은 (0,-1)(빈 점)에서 1을 지나 오르는 직선.
% 원문의 라벨 O·1·-1 만 둔다. 스캔 화질이 낮아 TikZ 로 다시 그림.
\begin{tikzpicture}[line width=0.55pt, x=0.6cm, y=0.6cm, every node/.style={font=\footnotesize}]
  \draw[->, >=stealth, line width=0.7pt] (-2.2,0) -- (2.3,0) node[below=1pt] {$x$};
  \draw[->, >=stealth, line width=0.7pt] (0,-1.8) -- (0,2.3) node[left=1pt] {$y$};
  \draw[line width=0.6pt] (-1.9,1.9) -- (0,0);
  \draw[line width=0.6pt] (0,-1) -- (1.95,0.95);
  \draw[fill=white, line width=0.5pt] (0,-1) circle (2.2pt);
  \fill (0,0) circle (2.2pt);
  \node[below left=0pt and -2pt] at (0,0) {$\pt{O}$};
  \node[below=1pt] at (1,-0.05) {$1$};
  \node[anchor=east, inner sep=2.5pt] at (-0.14,-1) {$-1$};
  \node[anchor=west] at (0.85,1.7) {$y=f(x)$};
\end{tikzpicture}
